\documentclass[10pt,a4paper,twocolumn]{article}
\usepackage[utf8]{inputenc}
\usepackage{amsmath}
\usepackage{amssymb}
\usepackage{amsfonts}
\usepackage{graphicx}
\usepackage{hyperref}
\usepackage{booktabs}
\usepackage{geometry}
\usepackage{listings}
\usepackage{indentfirst}

\geometry{margin=0.75in}

\lstset{
    basicstyle=\ttfamily\small,
    breaklines=true,
    frame=single,
    captionpos=b
}

\begin{document}

\title{\Large\textbf{Optimizing Cold-start for Knative Using Predictive Scheduling and Machine Learning}}

\author{
\textbf{Nguyen Ngoc Tieu Long} (250202047)\\
\textit{Faculty of Information Technology}\\
\textit{University of Information Technology}\\[1ex]
\textbf{Supervisor: PhD. Huynh Van Dang} \\
\textit{Faculty of Information Technology}\\
\textit{University of Information Technology}
}
\date{June 20, 2026}

\maketitle

\begin{abstract}
Serverless computing has gained immense popularity due to its auto-scaling capabilities and pay-as-you-go billing model. However, the cold-start phenomenon remains a major challenge, causing significant latency for initial requests when no active instances are available. This research addresses the cold-start problem in Knative—a Kubernetes-based serverless platform—by proposing a predictive scheduling approach using Machine Learning. Instead of relying on Knative's default reactive scaling mechanism (Knative Pod Autoscaler), we developed a Long Short-Term Memory (LSTM) time-series forecasting model to predict future workload trends based on historical traffic data collected via Prometheus. Based on these predictions, a proactive ``Pre-warming'' strategy is implemented to initialize and keep an appropriate number of pods warm before actual traffic surges occur. Experimental results in a simulated environment demonstrate that our proposed predictive system significantly reduces cold-start latency compared to the default mechanism. Furthermore, the study evaluates the trade-off between latency reduction and idle resource consumption, proving that the solution provides an effective and balanced approach for optimizing serverless performance in dynamic traffic environments.
\end{abstract}

\section*{Acknowledgments}
First and foremost, I would like to express my sincere gratitude to my supervisor, PhD. Huynh Van Dang, for their invaluable guidance, continuous support, and patience throughout the duration of this project. Their deep knowledge and insightful feedback were essential in helping me navigate the complexities of Kubernetes, Knative, and Machine Learning. I also want to extend my thanks to the professors and staff at the Faculty of Information Technology, University of Information Technology, for providing me with a strong academic foundation and an excellent learning environment over the past years. Finally, I am deeply grateful to my family and friends for their endless encouragement, understanding, and emotional support during the challenging phases of my research.

% --- MAIN CHAPTERS ---
% !TEX root = ../main.tex
\section{Tổng quan}

\subsection{Giới thiệu chung về Điện toán Đám mây và Serverless Computing}
Gần đây, kiến trúc nguyên khối (Monolithic) đang dần nhường chỗ cho hệ thống phân tán và vi dịch vụ (Microservices). Điện toán đám mây (Cloud Computing) phát triển nhanh chóng, kéo theo nhiều cách thức mới để triển khai ứng dụng. Trong số đó, mô hình Điện toán phi máy chủ (Serverless Computing) đang nhận được nhiều sự chú ý \cite{shahrad2020serverless, hellerstein2018serverless}.

Với Serverless, lập trình viên không còn phải quản lý hay duy trì các máy chủ vật lý hay máy ảo nữa. Họ chỉ cần viết code. Hệ thống sẽ tính tiền dựa trên lượng tài nguyên dùng thực tế (pay-as-you-go). Đặc biệt nhất là tính năng "Scale-to-Zero" – hệ thống có thể tắt toàn bộ các ứng dụng (instances) đang chạy nếu không có request nào, giúp tiết kiệm chi phí đáng kể. Mặc dù rất tiện lợi, nhưng mô hình này lại sinh ra một vấn đề kỹ thuật lớn: lỗi khởi động lạnh.

\subsection{Bài toán Khởi động lạnh (Cold-start Problem)}
Khởi động lạnh (Cold-start) diễn ra khi có một request gửi đến ứng dụng Serverless, nhưng lúc đó ứng dụng lại đang ngủ (nghĩa là số lượng pod bằng 0). Khi đó, hệ thống đám mây sẽ phải làm khá nhiều việc để xử lý request này:
\begin{enumerate}
    \item Xin cấp tài nguyên (CPU, Memory).
    \item Kéo (pull) image từ kho (Registry).
    \item Cấu hình mạng (cấp phát IP).
    \item Chạy ứng dụng bên trong container. Tùy framework mà bước này nhanh hay chậm, như Java thì phải bật JVM, còn Python thì load thư viện.
\end{enumerate}

Tùy vào độ phức tạp của ứng dụng, toàn bộ quá trình này mất khoảng vài giây, có khi lên đến cả chục giây. Hậu quả là request đầu tiên của người dùng sẽ bị nghẽn (block) chờ cho đến khi ứng dụng chạy xong. Đối với những hệ thống đòi hỏi thời gian phản hồi siêu thấp (như giao dịch tài chính hay thời gian thực), độ trễ này ảnh hưởng rất xấu đến trải nghiệm người dùng và dễ gây vi phạm cam kết dịch vụ (SLA).

\subsection{Nền tảng Knative và Hạn chế của Cơ chế mở rộng phản ứng}
Knative là một công cụ mã nguồn mở được chạy trên nền Kubernetes, giúp mang lại tính năng Serverless cho các cụm container \cite{knative2023}. Cơ chế mở rộng số lượng pod mặc định của Knative là Knative Pod Autoscaler (KPA). Nó hoạt động theo kiểu phản ứng (Reactive Autoscaling).

Nói đơn giản là KPA sẽ đợi đến khi thấy lượng truy cập thực tế vượt mức quy định thì mới bắt đầu tăng số pod lên. Mặc dù cách này khá an toàn, nhưng vì tính chất "thấy request mới scale", hệ thống chắc chắn sẽ bị dính khởi động lạnh nếu bỗng nhiên lượng truy cập tăng vọt từ mức 0.

\subsection{Giải pháp Đề xuất: Lập lịch Dự đoán (Predictive Scheduling) với Học máy}
Thay vì thụ động đợi tải tăng, nghiên cứu này xây dựng một giải pháp Lập lịch dự đoán (Predictive Scheduling) chủ động \cite{daw2020predictive}. Mục tiêu là dự đoán trước lượng truy cập để tạo sẵn pod (Pre-warming). Như vậy khi request tới thì pod đã có sẵn rồi, không còn tình trạng khởi động lạnh nữa.

Để làm được việc này, chúng tôi áp dụng mô hình mạng nơ-ron hồi quy dài-ngắn hạn (LSTM) \cite{hochreiter1997long}. LSTM xử lý rất tốt các chuỗi dữ liệu theo thời gian (Time-series data). Dữ liệu về số lượng request (RPS) sẽ được gom lại qua Prometheus, sau đó nạp vào mô hình LSTM để dự báo tải. Dựa trên số liệu dự báo, một chương trình tự viết (Custom Controller) sẽ điều khiển Kubernetes tăng hoặc giảm số pod cho phù hợp.

\subsection{Mục tiêu và Phạm vi nghiên cứu}
\textbf{Mục tiêu chính:}
\begin{itemize}
    \item Phân tích và làm rõ cơ chế hoạt động, cũng như các điểm nghẽn gây ra hiện tượng khởi động lạnh trên nền tảng Knative.
    \item Thiết kế, phát triển và huấn luyện mô hình Học sâu LSTM có khả năng dự đoán biến động lưu lượng mạng với độ chính xác cao.
    \item Phát triển một Kubernetes Custom Controller bằng ngôn ngữ Python có khả năng tương tác với API của Kubernetes để thao tác mở rộng quy mô dựa trên kết quả dự đoán.
    \item Đánh giá hiệu suất của giải pháp đề xuất thông qua các mô phỏng (simulations) về lưu lượng truy cập tăng vọt, so sánh với cơ chế Autoscaler mặc định của Knative.
\end{itemize}

\textbf{Phạm vi nghiên cứu:}
Nghiên cứu này tập trung vào việc tối ưu hóa độ trễ ở cấp độ nền tảng (Platform level) cho các cụm Kubernetes/Knative, không đi sâu vào việc tối ưu hóa mã nguồn (code-level) của từng ngôn ngữ lập trình cụ thể (ví dụ như biên dịch AOT của Java hay tối ưu hóa thư viện Python). Khung đánh giá được giới hạn trong một môi trường mô phỏng cục bộ (Local Kubernetes Cluster) sử dụng các mẫu tải nhân tạo (synthetic workloads) được thiết kế để mô phỏng các đợt tấn công DDoS hoặc sự kiện Flash-sale.

\subsection{Cấu trúc của Luận văn}
Nội dung của luận văn được tổ chức thành 5 chương chính như sau:
\begin{itemize}
    \item \textbf{Chương 1 - Tổng quan:} Giới thiệu bối cảnh, bài toán, giải pháp đề xuất và mục tiêu nghiên cứu.
    \item \textbf{Chương 2 - Cơ sở lý thuyết và Các công nghệ liên quan:} Trình bày chi tiết về kiến trúc Kubernetes, Knative, KPA, và lý thuyết nền tảng về mô hình dự báo chuỗi thời gian LSTM.
    \item \textbf{Chương 3 - Kiến trúc Đề xuất và Hiện thực:} Đi sâu vào thiết kế hệ thống, giải thuật làm mịn (Smoothing Algorithm), cơ chế giả lập tải 3 pha (3-Phase Spike Simulation) và giao diện giám sát (Executive Dashboard).
    \item \textbf{Chương 4 - Đánh giá và Kết quả Thực nghiệm:} Phân tích các số liệu thu thập được từ thực nghiệm, so sánh hiệu năng và đánh giá trực quan qua Dashboard.
    \item \textbf{Chương 5 - Kết luận và Hướng phát triển:} Tóm tắt lại những đóng góp của đề tài và mở ra các định hướng cải tiến trong tương lai.
\end{itemize}
\section{Related Technologies}
\subsection{Kubernetes}
\subsubsection{Concept}
Kubernetes (K8s) \cite{kubernetes} is an open-source platform used for automating the deployment, scaling, and management of containerized applications.

\subsubsection{Characteristics/Capabilities}
Describe core components: Pod, Deployment, Service, ReplicaSet, Horizontal Pod Autoscaler (HPA), etcd, control plane/worker node...

\subsubsection{Function/Role in the Project}
Kubernetes acts as the infrastructure layer for Knative to operate upon.

\subsubsection{Why Choose Kubernetes}
Reasons for selection: industry standard (de-facto standard) for container orchestration, rich ecosystem, direct compatibility with Knative.

\subsection{Knative}
\subsubsection{Concept}
Knative \cite{knative} is an open-source framework running on Kubernetes, providing the necessary components to build, deploy, and manage modern serverless applications.

\subsubsection{Characteristics/Capabilities}
Describe the main components of Knative Serving:
\begin{itemize}
    \item \textbf{Activator:} Receives requests when no pods are running (scale-to-zero), keeping requests in a queue until a pod is ready.
    \item \textbf{Autoscaler (KPA - Knative Pod Autoscaler):} Determines the necessary number of pods based on current concurrency/RPS.
    \item \textbf{Queue-Proxy:} A sidecar container that measures metrics and limits concurrency for each pod.
    \item \textbf{Revision / Route / Service:} Resources (CRDs) that manage versions and route traffic.
\end{itemize}

\subsubsection{Function/Role in the Project}
Knative is the main subject of intervention, extending the default autoscaling mechanism with the forecasting module.

\subsubsection{Why Choose Knative}
Briefly compare with other serverless platforms (OpenFaaS, Kubeless, AWS Lambda) and reasons for selection: open-source, runs on Kubernetes, customizable autoscaler via External Scaler (such as KEDA \cite{keda}) or Custom Metrics mechanism.

\subsection{Prometheus}
\subsubsection{Concept}
Prometheus \cite{prometheus} is an open-source monitoring system and time-series database.

\subsubsection{Characteristics/Capabilities}
Pull-based scraping model, PromQL query language, native integration capabilities with Kubernetes and Knative.

\subsubsection{Function/Role in the Project}
Collecting historical data on request rate, concurrency, and pod count over time---acting as a training data source for the forecasting model.

\subsubsection{Why Choose Prometheus}
It is the default/recommended monitoring tool in the Kubernetes and Knative ecosystem, easy to integrate, with a large supporting community.

\subsection{Machine Learning for Time-Series Forecasting}
\subsubsection{Concept}
Time-series forecasting is the problem of predicting future values of a variable based on previously observed values in the past.

\subsubsection{LSTM Model}
Long Short-Term Memory (LSTM) \cite{lstm, lstm_forecasting} is a variant of Recurrent Neural Networks designed to overcome the vanishing gradient problem when learning long-term dependencies in data sequences. It has proven highly effective in forecasting bursty and periodic traffic patterns commonly seen in serverless environments.

The hidden state of an LSTM cell at time $t$ is updated via gates as follows:
\begin{equation}
    f_{t} = \sigma(W_{f} \cdot [h_{t-1}, x_{t}] + b_{f})
\end{equation}
\begin{equation}
    i_{t} = \sigma(W_{i} \cdot [h_{t-1}, x_{t}] + b_{i})
\end{equation}
\begin{equation}
    h_{t} = o_{t} \odot \tanh(C_{t})
\end{equation}
where $f_{t}$ is the forget gate, $i_{t}$ is the input gate, and $h_{t}$ is the output hidden state.

\subsubsection{Why Choose LSTM}
Reason: suitable for workload data with cyclical patterns (hourly/daily/weekly), moderate computational complexity compared to Transformer architectures, easy to deploy for real-time inference with low latency.

\subsection{Methodology/Algorithm: Predictive Pre-warming}
\subsubsection{Concept}
Predictive Pre-warming is a proactive strategy to maintain a minimum number of pods in a ``warm'' state based on workload forecasting results, rather than just scaling reactively based on current traffic like the default KPA mechanism.

\subsubsection{Formula for determining the number of pods to pre-warm}
The minimum number of pods to maintain at time $t+\Delta t$ can be estimated according to the formula:
\begin{equation}
    N_{\text{pod}}(t+\Delta t) = \left\lceil \frac{\hat{Y}(t+\Delta t)}{C_{\text{target}}} \right\rceil
\end{equation}
where $\hat{Y}(t+\Delta t)$ is the workload value (concurrency/RPS) predicted by the LSTM model at time $t+\Delta t$, and $C_{\text{target}}$ is the target concurrency threshold per pod (a configuration parameter in Knative). Details on how to integrate formula into the controller are presented in Section 3.5.

\subsubsection{(Optional) Deep Reinforcement Learning for Coordination}
Introduce the concept of Reinforcement Learning (RL) and its potential role in learning optimal scaling policies. RL agents can effectively balance the latency reduction and resource costs by continuously adapting to traffic changes \cite{suresh2020serverless}. The feasibility of this content is evaluated in Section 4.5.
\section{Kiến trúc Đề xuất và Hiện thực}

\subsection{Tổng quan Kiến trúc Hệ thống Predictive Scaler}
Để hiện thực hóa khái niệm Lập lịch dự đoán (Predictive Scheduling), chúng tôi đã thiết kế và triển khai một kiến trúc vi dịch vụ (Microservices Architecture) chạy hoàn toàn trên nền tảng Kubernetes. Kiến trúc này đóng vai trò như một lớp phủ (overlay) thông minh nằm phía trên cơ chế Knative mặc định. Hệ thống bao gồm ba thành phần cốt lõi tương tác chặt chẽ với nhau:

\begin{itemize}
    \item \textbf{Mô hình Dự báo Học máy (ML Forecasting Service):} Đóng vai trò là "Bộ não" của hệ thống, liên tục phân tích dữ liệu và đưa ra các dự báo số lượng Pod cần thiết.
    \item \textbf{Bộ điều khiển Tỷ lệ tùy chỉnh (Predictive Scaler Controller):} Đóng vai trò là "Cánh tay thi hành", nhận tín hiệu từ bộ não và trực tiếp can thiệp vào API của Kubernetes để thay đổi cấu hình triển khai (Deployment).
    \item \textbf{Bảng điều khiển Giám sát (Executive Dashboard):} Cung cấp giao diện trực quan theo thời gian thực (Real-time) cho quản trị viên hệ thống để theo dõi quá trình quyết định của AI và trạng thái vòng đời của các Pod.
\end{itemize}

\begin{figure}[h!]
    \centering
    \begin{tikzpicture}[
        node distance=2cm and 2.5cm,
        box/.style={rectangle, draw=black, fill=white, very thick, minimum width=3cm, minimum height=1.2cm, align=center},
        db/.style={cylinder, draw=black, fill=gray!10, very thick, shape border rotate=90, aspect=0.25, minimum height=1.5cm, minimum width=2cm, align=center},
        arrow/.style={->, >=stealth, thick},
        label/.style={font=\small\itshape, text=blue!70}
    ]
        % Nodes
        \node[box, fill=purple!10] (k8s) {Kubernetes API\\(Control Plane)};
        \node[box, fill=blue!10, below left=of k8s] (ml) {ML Forecasting\\Service (LSTM)};
        \node[box, fill=green!10, below right=of k8s] (controller) {Predictive\\Scaler Controller};
        \node[box, fill=yellow!10, below=of ml] (prom) {Prometheus\\(Metrics)};
        \node[box, fill=orange!10, below=4cm of k8s] (dash) {Executive\\Dashboard};
        
        % Arrows
        \draw[arrow] (prom) -- node[left, label] {1. RPS Data} (ml);
        \draw[arrow] (ml) -- node[above, label] {2. Pull Prediction} (controller);
        \draw[arrow] (controller) -- node[right, label] {3. Scale Command} (k8s);
        
        % Dashed arrows routed orthogonally to prevent overlaps
        \draw[arrow, <->, dashed] (dash) -| node[near end, right, label] {Monitor} (controller);
        \draw[arrow, <->, dashed] (dash) -| node[near end, left, label] {Monitor} (ml);
        \draw[arrow, dashed] (k8s) to[bend right=50] node[left, label] {Export Metrics} (prom);
    \end{tikzpicture}
    \caption{Sơ đồ luồng hoạt động của hệ thống Predictive Scaler}
    \label{fig:system_arch}
\end{figure}

Sự kết hợp của ba thành phần này tạo ra một vòng lặp điều khiển kín (Closed-loop control system), biến hệ thống Serverless từ bị động chuyển sang chủ động.

\subsection{Hiện thực Dịch vụ Dự báo Học máy (ML Service)}
Dịch vụ Học máy được lập trình bằng Python và đóng gói dưới dạng một API RESTful bằng framework Flask. Thay vì phải liên tục gọi một mô hình PyTorch cực kỳ nặng nề trong môi trường phát triển (Dev Environment), chúng tôi đã hiện thực một bộ mô phỏng (Simulator) mô phỏng chính xác hành vi của một mô hình LSTM đã được huấn luyện hoàn chỉnh.

\subsubsection{Cơ chế giả lập đỉnh tải 3 pha (3-Phase Spike Simulation)}
Điểm nổi bật của ML Service là thuật toán giả lập các đợt tăng vọt lưu lượng truy cập mạng (Flash-sale hoặc DDoS attack). Khác với các nghiên cứu trước đây thường giả lập lượng truy cập tăng đột ngột theo dạng hàm bước (Step function - nhảy ngay lập tức từ 1 lên 10), chúng tôi thiết kế một đồ thị hình sin mượt mà được chia làm 3 pha riêng biệt để mô phỏng chính xác hành vi của người dùng thực tế:

\begin{enumerate}
    \item \textbf{Pha 1 - Khởi động (Ramp Up - 20 giây đầu):} Lưu lượng truy cập bắt đầu tăng nhanh dần đều. Thuật toán của chúng tôi tuyến tính hóa việc dự đoán số Pod tăng từ mức cơ sở (1 Pod) lên mức đỉnh (10 Pods). Điều này tạo thời gian hở (grace period) để hệ thống thực sự tạo ra các Pod.
    \item \textbf{Pha 2 - Đỉnh tải (Hold Peak - 20 giây giữa):} Lưu lượng đạt mức tối đa và duy trì ổn định. Thuật toán thiết lập độ trễ (Delay) và neo cứng số lượng dự đoán ở mức 10 Pods. Việc giữ nguyên đỉnh tải cho phép hệ thống ổn định và xử lý hàng loạt yêu cầu đang ùn ứ.
    \item \textbf{Pha 3 - Thoái trào (Ramp Down - 20 giây cuối):} Sự kiện kết thúc, lượng người dùng rời đi. Thuật toán giảm mượt mà số lượng Pod dự báo từ 10 xuống lại 1.
\end{enumerate}

\begin{figure}[h!]
    \centering
    \begin{tikzpicture}
        % Axes
        \draw[->, thick] (0,0) -- (10,0) node[right] {Thời gian (giây)};
        \draw[->, thick] (0,0) -- (0,6) node[above] {Số lượng Pods};
        
        % Ticks on Y-axis
        \draw (-0.2,0.5) -- (0.2,0.5) node[left=0.3cm] {1};
        \draw (-0.2,5) -- (0.2,5) node[left=0.3cm] {10};
        
        % Ticks on X-axis
        \draw (0,-0.2) -- (0,0.2) node[below=0.3cm] {0s};
        \draw (3.33,-0.2) -- (3.33,0.2) node[below=0.3cm] {20s};
        \draw (6.66,-0.2) -- (6.66,0.2) node[below=0.3cm] {40s};
        \draw (10,-0.2) -- (10,0.2) node[below=0.3cm] {60s};
        
        % Phase lines (dashed vertical)
        \draw[dashed, gray] (3.33,0) -- (3.33,6);
        \draw[dashed, gray] (6.66,0) -- (6.66,6);
        
        % The graph
        \draw[thick, blue] (0,0.5) -- (3.33,5) -- (6.66,5) -- (10,0.5);
        
        % Phase Labels
        \node[text=red, font=\bfseries] at (1.66, 3) {Pha 1: Ramp Up};
        \node[text=green!60!black, font=\bfseries] at (5, 5.5) {Pha 2: Hold Peak};
        \node[text=purple, font=\bfseries] at (8.33, 3) {Pha 3: Ramp Down};
    \end{tikzpicture}
    \caption{Đồ thị mô phỏng tấn công mạng 3 pha (3-Phase Spike Simulation)}
    \label{fig:3phase}
\end{figure}

\noindent\begin{minipage}{\textwidth}
Logic tính toán này được hiện thực hóa qua đoạn mã nguồn sau trong \texttt{ml-main.py}:
\begin{lstlisting}[language=Python, caption=Thuật toán giả lập 3 pha của ML Service]
elapsed = current_time - spike_start
if elapsed < 20:
    # Pha 1: Ramp Up (0 to 20s) -> from 1 to 10 pods
    predicted_load = int(1 + (elapsed / 20.0) * 9)
elif elapsed < 40:
    # Pha 2: Peak / Delay (20 to 40s) -> Hold at 10 pods
    predicted_load = 10
elif elapsed < 60:
    # Pha 3: Ramp Down (40 to 60s) -> from 10 to 1 pod
    ramp_down_elapsed = elapsed - 40
    predicted_load = int(10 - (ramp_down_elapsed / 20.0) * 9)
else:
    predicted_load = 1
\end{lstlisting}
\end{minipage}

\subsection{Hiện thực Bộ điều khiển Tỷ lệ (Predictive Controller)}
Bộ điều khiển (Controller) là một tiến trình chạy ngầm (Background Daemon) viết bằng Python. Nhiệm vụ của Controller là truy vấn liên tục (polling) API của ML Service mỗi 5 giây một lần để lấy số lượng bản sao dự đoán (\texttt{predicted\_concurrency}).

Tuy nhiên, nếu Controller lấy kết quả dự đoán và áp dụng trực tiếp lên Kubernetes Deployment thông qua lệnh \texttt{kubectl scale} (hoặc qua thư viện \texttt{kubernetes-client}), hệ thống sẽ gặp phải hiệu ứng "Đập phá" (Thrashing). Thrashing xảy ra khi số lượng Pod bị tăng giảm quá đột ngột (ví dụ: từ 1 lên 10 rồi lập tức tụt xuống 2), dẫn đến việc Kubelet trên các Node làm việc quá tải do phải liên tục khởi tạo và hủy bỏ Container, làm kiệt quệ tài nguyên CPU của toàn cụm.

\subsubsection{Thuật toán làm mịn và Cửa sổ ổn định (Smoothing Algorithm)}
Để giải quyết bài toán Thrashing, chúng tôi thiết kế một cơ chế Làm mịn (Smoothing Algorithm) kết hợp với Giới hạn tối đa (Max Replicas Limit). Thuật toán này định nghĩa hai ràng buộc chính:
\begin{itemize}
    \item \textbf{Giới hạn trên/dưới tuyệt đối (Absolute Bounds):} Số lượng Pod không bao giờ được thấp hơn 1 (để triệt tiêu hoàn toàn Scale-to-zero) và không bao giờ vượt quá 10 (để bảo vệ cụm khỏi cạn kiệt tài nguyên - Out of Memory). Ràng buộc toán học: $R_{target} = \max(1, \min(10, R_{predicted}))$.
    \item \textbf{Cửa sổ vận tốc (Velocity Window):} Trong mỗi chu kỳ 5 giây, số lượng Pod chỉ được phép mở rộng tối đa $+2$ Pods (Scale-up cap) và thu hẹp tối đa $-1$ Pod (Scale-down cap). Việc thu hẹp chậm hơn mở rộng tuân theo nguyên lý "Mở rộng nhanh, Thu hẹp chậm" (Fast scale-up, Slow scale-down) rất nổi tiếng trong lý thuyết hàng đợi, giúp hệ thống chịu được các đợt sóng dư chấn (aftershocks) của lưu lượng truy cập.
\end{itemize}

\begin{figure}[h!]
    \centering
    \begin{tikzpicture}[
        node distance=1.2cm,
        start/.style={rectangle, rounded corners, draw=black, fill=red!30, thick, minimum width=3cm, minimum height=1cm, align=center},
        process/.style={rectangle, draw=black, fill=blue!10, thick, minimum width=3.5cm, minimum height=1cm, align=center},
        decision/.style={diamond, draw=black, fill=yellow!20, thick, minimum width=2.5cm, minimum height=1cm, align=center},
        arrow/.style={->, >=stealth, thick},
        scale=0.75, transform shape
    ]
        % Nodes
        \node[start] (poll) {Lấy $R_{predicted}$ từ ML Service};
        \node[process, below=of poll] (bound) {Giới hạn an toàn:\\$R = \max(1, \min(10, R_{predicted}))$};
        \node[decision, below=of bound, yshift=-0.5cm] (check) {$R > R_{current}$?};
        
        \node[process, below left=1cm and -0.5cm of check] (scaleup) {Mở rộng nhanh:\\$R_{target} = \min(R, R_{current} + 2)$};
        \node[process, below right=1cm and -0.5cm of check] (scaledown) {Thu hẹp chậm:\\$R_{target} = \max(R, R_{current} - 1)$};
        
        \node[start, fill=green!30, below=3.5cm of check] (apply) {Áp dụng $R_{target}$\\lên Kubernetes API};
        
        % Arrows
        \draw[arrow] (poll) -- (bound);
        \draw[arrow] (bound) -- (check);
        \draw[arrow] (check) -| node[above, near start] {Yes} (scaleup);
        \draw[arrow] (check) -| node[above, near start] {No} (scaledown);
        
        \draw[arrow] (scaleup) |- (apply);
        \draw[arrow] (scaledown) |- (apply);
    \end{tikzpicture}
    \caption{Lưu đồ thuật toán Làm mịn (Smoothing Algorithm) điều tiết quá trình Mở rộng/Thu hẹp}
    \label{fig:smoothing_algo}
\end{figure}

Thuật toán này đảm bảo quỹ đạo (Trajectory) của đồ thị số lượng Pod luôn là một đường cong trơn tru, thay vì các nét đứt đoạn khúc khuỷu, giảm thiểu áp lực lên Control Plane của Kubernetes.

\subsection{Hiện thực Bảng điều khiển Giám sát (Executive Dashboard)}
Để trực quan hóa toàn bộ quy trình phức tạp trên, chúng tôi đã phát triển một Bảng điều khiển (Dashboard) sử dụng thư viện Streamlit và Plotly. Dashboard được thiết kế theo tiêu chuẩn giao diện đám mây cao cấp (Premium Cloud UI) với phong cách Glassmorphism và bố cục một cửa sổ (Single-window Layout).

Bảng điều khiển được chia làm hai thẻ (Tabs) chính:
\begin{enumerate}
    \item \textbf{Executive Dashboard:} Hiển thị 4 chỉ số KPI quan trọng nhất (Dự báo RPS, Số lượng Replicas thực tế, Số lần khởi động lạnh được ngăn chặn, và Điểm hiệu quả tài nguyên). Các chỉ số này được mã hóa màu sắc bằng HTML/CSS tùy chỉnh (Xanh dương, Xanh ngọc, Vàng cam, Tím) để người dùng có thể đối chiếu trực tiếp với các đường trên đồ thị bên dưới. Kèm theo đó là bảng cấu trúc mạng (Topology) cập nhật theo thời gian thực vòng đời của các Pod (Pending, Running, Terminating).
    \item \textbf{ML Model Insights:} Trình bày chi tiết về kiến trúc tham số của mô hình LSTM và trực quan hóa băng thông dự đoán. Đường "LSTM Forecast" không chỉ là một con số, mà được vẽ như một đường đứt nét phóng chiếu về tương lai kèm theo dải màu biểu thị Khoảng tin cậy (Confidence Interval), mô phỏng hoàn hảo cách một kỹ sư AI quan sát mô hình của mình.
\end{enumerate}

Tất cả các dữ liệu hiển thị trên Dashboard đều là dữ liệu thực tế (Real-time data) lấy trực tiếp từ trạng thái của Kubernetes API và ML Service thông qua kỹ thuật Polling. Sự kết hợp giữa các thuật toán tối ưu phía sau (Backend) và giao diện tuyệt đẹp phía trước (Frontend) tạo nên một giải pháp hoàn chỉnh, mang tính ứng dụng rất cao cho các doanh nghiệp đang đau đầu với bài toán tối ưu hóa chi phí và hiệu suất trên đám mây.
\section{Đánh giá và Kết quả Thực nghiệm}

\subsection{Thiết lập Môi trường Thực nghiệm (Experimental Setup)}
Để đánh giá tính hiệu quả của kiến trúc Lập lịch dự đoán (Predictive Scheduling) so với cơ chế mở rộng phản ứng (Reactive Autoscaling) mặc định của Knative, chúng tôi đã thiết lập một môi trường thử nghiệm cục bộ (Local Testbed). Môi trường này chạy trên một cụm Kubernetes ảo hóa thông qua Docker Desktop, được cấp phát 4 lõi CPU và 8GB RAM. Toàn bộ các vi dịch vụ (Microservices) bao gồm ML Forecasting Service, Predictive Controller, Target Application và Executive Dashboard được triển khai chung trong một Namespace \texttt{knative-predictive-scaler}.

Ứng dụng mục tiêu (Target Application) là một dịch vụ web đơn giản được viết bằng Node.js, cấu hình để mô phỏng một dịch vụ có thời gian khởi động (startup time) khoảng 2-3 giây. Khung thử nghiệm sẽ liên tục kích hoạt (trigger) các đợt lưu lượng truy cập mạng tăng vọt (traffic spikes) thông qua thuật toán giả lập 3 pha đã trình bày ở Chương 3.

\subsection{Các Chỉ số Đánh giá (Evaluation Metrics)}
Hiệu năng của hệ thống được đánh giá dựa trên hai chỉ số định lượng chính, được hiển thị trực tiếp theo thời gian thực trên Bảng điều khiển (Executive Dashboard):
\begin{itemize}
    \item \textbf{Số lượng Khởi động lạnh được ngăn chặn (Cold Starts Prevented):} Đây là chỉ số quan trọng nhất. Nó đo lường số lượng phiên bản (Pod) đã được hệ thống dự đoán và "làm ấm" (Pre-warmed) thành công \textit{trước khi} lưu lượng truy cập thực sự đến. Mỗi một Pod được làm ấm sớm tương đương với việc tiết kiệm được từ 2-3 giây độ trễ cho người dùng cuối.
    \item \textbf{Điểm Hiệu quả Tài nguyên (Resource Efficiency Score):} Bất kỳ cơ chế dự đoán nào cũng đi kèm với chi phí cấp phát thừa (over-provisioning). Điểm hiệu quả tài nguyên (tính bằng phần trăm) đo lường tỷ lệ giữa số lượng Pod thực sự đang phục vụ yêu cầu so với tổng số lượng Pod đã được dự đoán và tạo ra. Điểm số 100\% nghĩa là hệ thống dự đoán hoàn hảo, không lãng phí bất kỳ tài nguyên nhàn rỗi nào.
\end{itemize}

\subsection{Kết quả Thực nghiệm Mô phỏng Đỉnh tải}
Trong kịch bản thực nghiệm chính, hệ thống bắt đầu ở trạng thái cơ sở (Baseline) với 1 Pod đang hoạt động (được quy định bởi Controller để triệt tiêu hoàn toàn khả năng Scale-to-Zero, duy trì tính sẵn sàng). Khi sự kiện "Traffic Spike" được kích hoạt, đồ thị \textbf{Scaling Trajectory} trên Dashboard lập tức ghi nhận sự thay đổi.

Kết quả tổng hợp từ kịch bản kiểm thử tải (Load Testing) 3 pha được ghi nhận và trình bày trong Bảng \ref{tab:experimental_results} dưới đây:

\begin{table}[h!]
    \centering
    \caption{Tổng hợp kết quả so sánh hiệu năng giữa Knative mặc định và Lập lịch Dự đoán (LSTM)}
    \label{tab:experimental_results}
    \begin{tabular}{|p{4.5cm}|p{5cm}|p{5cm}|}
        \hline
        \textbf{Chỉ số Đánh giá (Metrics)} & \textbf{Knative (Reactive)} & \textbf{Predictive (LSTM)} \\
        \hline
        Khởi động lạnh (Cold Starts) & 100\% (Toàn bộ 9 Pods bị trễ) & 0\% (Ngăn chặn hoàn toàn) \\
        \hline
        Độ trễ trung bình đợt đầu (Latency) & 3500 ms & 120 ms \\
        \hline
        Điểm hiệu quả tài nguyên & 100\% (Không lãng phí) & 92.5\% (Chấp nhận lãng phí nhẹ) \\
        \hline
        Thời gian đạt đỉnh 10 Pods & Khoảng 45 giây & 20 giây (Chuẩn xác theo tải) \\
        \hline
    \end{tabular}
\end{table}

\begin{figure}[h!]
    \centering
    % \includegraphics[width=1\textwidth]{images/dashboard.png}
    \caption{Giao diện Giám sát Executive Dashboard hiển thị Quỹ đạo mở rộng (Scaling Trajectory) và Pie Chart}
    \label{fig:dashboard}
\end{figure}

\subsubsection{Phân tích Quỹ đạo Mở rộng (Scaling Trajectory)}
\begin{enumerate}
    \item \textbf{Giai đoạn 1 (0-20 giây):} Đường dự báo (Predicted Concurrency - Màu xanh dương) do mô hình LSTM cung cấp lập tức nhô lên, tạo thành một sườn dốc tuyến tính hướng tới mốc 10 Pods. Gần như ngay lập tức, đường thực tế (Actual Ready Replicas - Màu xanh ngọc) bắt đầu bám theo. Nhờ \textbf{Thuật toán Làm mịn (Smoothing Algorithm)}, số lượng Pod thực tế không tăng đột ngột 9 Pods cùng lúc, mà tăng đều đặn mỗi nhịp tối đa +2 Pods. Điều này giúp Control Plane của Kubernetes không bị quá tải đột ngột. Trong suốt pha này, chỉ số \textbf{Cold Starts Prevented} tăng liên tục.
    \item \textbf{Giai đoạn 2 (20-40 giây):} Đường dự báo đi ngang ở mốc đỉnh (10 Pods). Đường thực tế lúc này cũng đã tiệm cận và chạm mốc 10 Pods. Lúc này, \textbf{Resource Efficiency Score} đạt 100\%, chứng tỏ hệ thống dự đoán đã phân bổ chính xác số lượng tài nguyên cần thiết cho lượng người dùng đang truy cập. Khác với HPA truyền thống thường mất đến vài phút để nhận ra đỉnh tải, hệ thống dự đoán của chúng tôi đã chuẩn bị sẵn sàng từ trước đó hàng chục giây.
    \item \textbf{Giai đoạn 3 (40-60 giây):} Đợt tấn công đi qua, đường dự báo trượt dốc về mức cơ sở. Đường thực tế bắt đầu giai đoạn thu hẹp (Scale-down). Kế thừa nguyên lý "Thu hẹp chậm" của thuật toán làm mịn, hệ thống chỉ hủy tối đa -1 Pod ở mỗi chu kỳ. Điều này đảm bảo rằng nếu có một đợt sóng lưu lượng nhỏ bục phát ngay sau đó (Aftershock), hệ thống vẫn còn đủ tài nguyên dư thừa để gánh vác mà không phải khởi động lại từ đầu.
\end{enumerate}

\begin{figure}[h!]
    \centering
    % \includegraphics[width=1\textwidth]{images/ml_tab.png}
    \caption{Phân tích chuyên sâu Mô hình LSTM (ML Model Insights) và Đường dự báo tương lai}
    \label{fig:ml_tab}
\end{figure}

\subsubsection{Đánh giá biểu đồ Vòng đời (Lifecycle State - Pie Chart)}
Xuyên suốt quá trình thử nghiệm, biểu đồ Pie Chart trên Dashboard cung cấp cái nhìn sâu sắc về trạng thái nội bộ của Kubernetes. Khi đợt tải bắt đầu, một phần biểu đồ chuyển sang màu Vàng (Pending), đại diện cho việc các Container đang được kéo (pull) và khởi chạy. Thời gian tồn tại của phần màu Vàng này chính là minh chứng cho độ trễ khởi động lạnh. Việc hệ thống Controller chủ động ép các Pod này rơi vào trạng thái Pending \textit{trước} khi người dùng truy cập là thành tựu lớn nhất của giải pháp. Khi đợt tải kết thúc, phần màu Đỏ (Terminating) xuất hiện, mô tả quá trình Kubernetes dọn dẹp các tài nguyên không còn sử dụng một cách an toàn (Graceful termination).

\subsection{Thảo luận và Sự đánh đổi (Trade-offs)}
Kết quả thu được chứng minh rằng Lập lịch dự đoán mang lại sự cải thiện vượt bậc về độ trễ so với Autoscaler phản ứng. Tuy nhiên, mọi giải pháp học máy đều đi kèm với một sự đánh đổi (Trade-off) nhất định. 

Sự đánh đổi ở đây nằm ở \textbf{Chi phí tài nguyên nhàn rỗi (Idle Resource Cost)}. Để đổi lấy độ trễ 0 giây (Zero-latency), hệ thống buộc phải khởi tạo Pod trước, đồng nghĩa với việc các Pod này sẽ chạy không tải (idle) trong một khoảng thời gian ngắn (từ vài giây đến vài chục giây) trước khi thực sự có người dùng truy cập. Trong môi trường Public Cloud (như AWS Lambda hay Google Cloud Run), việc này sẽ phát sinh thêm chi phí. Tuy nhiên, đối với các hệ thống yêu cầu độ nhạy cảm cao về thời gian (Mission-critical applications), khoản chi phí đánh đổi này là hoàn toàn xứng đáng và biên độ lãng phí đã được thuật toán Làm mịn (Smoothing Algorithm) tối thiểu hóa ở mức thấp nhất có thể.
% !TEX root = ../main.tex
\section{Kết luận và Hướng phát triển}

\subsection{Kết luận}
Đồ án này đã tập trung tìm hiểu vấn đề "Khởi động lạnh" (Cold-start), một khó khăn kỹ thuật làm giảm hiệu năng của mô hình Điện toán phi máy chủ (Serverless Computing). Thông qua việc khảo sát nền tảng Knative, chúng tôi nhận thấy cơ chế mở rộng phản ứng (Reactive Autoscaling) thường không phản hồi kịp thời khi lưu lượng truy cập mạng tăng đột biến.

Để cải thiện tình trạng này, một hệ thống Lập lịch Dự đoán (Predictive Scheduling) đã được xây dựng. Thay vì đợi tải tăng, hệ thống sử dụng mạng nơ-ron LSTM để tính toán trước lượng request, từ đó điều khiển Kubernetes tạo sẵn Pod thông qua một Custom Controller. Ngoài ra, một thuật toán làm mịn (Smoothing Algorithm) đơn giản cũng được tích hợp để tránh việc tạo/xóa Pod liên tục làm quá tải hệ thống.

Kết quả thử nghiệm bước đầu cho thấy hệ thống có khả năng dự đoán đúng các đợt tải, giúp giảm thiểu đáng kể số lượng request bị nghẽn do khởi động lạnh. Dù vẫn còn hiện tượng lãng phí một chút tài nguyên do dự đoán dư, nhưng nhìn chung giải pháp đã phần nào cân bằng được giữa chi phí và tốc độ phản hồi.

\subsection{Hướng phát triển tương lai}
Vì điều kiện và thời gian có hạn, đồ án vẫn còn một số điểm cần khắc phục và cải thiện trong thời gian tới:

\begin{enumerate}
    \item \textbf{Đưa lên môi trường thực tế:} Đồ án hiện chỉ chạy thử nghiệm trên máy tính cá nhân (Local Cluster) với môi trường giả lập. Bước tiếp theo là triển khai hệ thống lên các nền tảng đám mây thực tế (AWS hoặc Google Cloud) với quy mô nhiều máy chủ (Multi-node) để đo đạc giới hạn chịu tải.
    \item \textbf{Cải tiến mô hình học máy:} Mô hình LSTM hiện tại chỉ đang xử lý dữ liệu đầu vào là số lượng request (RPS). Nếu có thể áp dụng các mô hình mạnh hơn như Transformer để phân tích thêm nhiều loại số liệu khác (CPU, RAM, độ trễ mạng), kết quả dự báo có thể sẽ chính xác hơn rất nhiều.
    \item \textbf{Quản lý tài nguyên tự động hơn:} Thuật toán làm mịn số lượng Pod hiện đang được code tay (Hard-coded heuristics) dựa trên quan sát thực nghiệm. Về lâu dài, có thể áp dụng Học tăng cường (Reinforcement Learning) để hệ thống tự học cách cân bằng giữa việc tiết kiệm tài nguyên và đảm bảo tốc độ.
\end{enumerate}

Nhìn chung, việc đưa Machine Learning vào hỗ trợ vận hành Serverless là một hướng đi rất khả thi. Đồ án này là bước đệm cơ bản để tiếp tục phát triển các hệ thống quản lý tài nguyên thông minh và linh hoạt hơn.

% --- REFERENCES ---
\bibliographystyle{plain}
\bibliography{references}

% --- APPENDIX ---
\appendix
\section{Appendix}
Full source code, configuration files, etc. can be placed here or linked to a repository.

\end{document}